% ═══════════════════════════════════════════════════════════════════════════
% FehrAdvice Executive Summary: US-Iran Crisis Model
% Created: 2026-01-28
% Style: FehrAdvice Corporate (CLAUDE.md REF-STYLE)
% ═══════════════════════════════════════════════════════════════════════════

\documentclass[11pt,a4paper]{article}

% ───────────────────────────────────────────────────────────────────────────
% PACKAGES
% ───────────────────────────────────────────────────────────────────────────
\usepackage[utf8]{inputenc}
\usepackage[T1]{fontenc}
\usepackage[ngerman]{babel}
\usepackage{geometry}
\usepackage{xcolor}
\usepackage{tikz}
\usepackage{graphicx}
\usepackage{fontspec}
\usepackage{tabularx}
\usepackage{booktabs}
\usepackage{fancyhdr}
\usepackage{lastpage}
\usepackage{hyperref}
\usepackage{parskip}
\usepackage{enumitem}

% ───────────────────────────────────────────────────────────────────────────
% FEHRADVICE COLORS (from CLAUDE.md)
% ───────────────────────────────────────────────────────────────────────────
\definecolor{fadarkblue}{RGB}{2,64,121}      % #024079 - Headlines, Links
\definecolor{falightblue}{RGB}{84,158,222}   % #549EDE - Secondary
\definecolor{fadarkgray}{RGB}{37,33,42}      % #25212A - Body text
\definecolor{falightgray}{RGB}{243,245,247}  % #F3F5F7 - Backgrounds
\definecolor{faorange}{RGB}{222,157,62}      % Accent
\definecolor{fared}{RGB}{192,57,43}          % Warning/Alert
\definecolor{fagreen}{RGB}{39,174,96}        % Success

% ───────────────────────────────────────────────────────────────────────────
% PAGE SETUP
% ───────────────────────────────────────────────────────────────────────────
\geometry{
    a4paper,
    top=20mm,
    bottom=20mm,
    left=20mm,
    right=20mm
}

% ───────────────────────────────────────────────────────────────────────────
% FONTS (FehrAdvice Style)
% ───────────────────────────────────────────────────────────────────────────
\setmainfont{DejaVu Sans}
\setsansfont{DejaVu Sans}

% ───────────────────────────────────────────────────────────────────────────
% HEADER & FOOTER
% ───────────────────────────────────────────────────────────────────────────
\pagestyle{fancy}
\fancyhf{}
\renewcommand{\headrulewidth}{0pt}
\fancyfoot[L]{\small\color{fadarkgray} FehrAdvice \& Partners AG | Evidence-Based Framework}
\fancyfoot[R]{\small\color{fadarkgray} Seite \thepage}

% ───────────────────────────────────────────────────────────────────────────
% CUSTOM COMMANDS
% ───────────────────────────────────────────────────────────────────────────
\newcommand{\sectionbox}[2]{%
    \noindent\colorbox{fadarkblue}{%
        \parbox{\dimexpr\textwidth-2\fboxsep}{%
            \color{white}\textbf{\sffamily #1}%
        }%
    }%
    \vspace{2mm}

    #2
    \vspace{4mm}
}

\newcommand{\highlightbox}[1]{%
    \noindent\fcolorbox{fadarkblue}{falightgray}{%
        \parbox{\dimexpr\textwidth-2\fboxsep-2\fboxrule}{%
            #1%
        }%
    }%
    \vspace{3mm}
}

% ───────────────────────────────────────────────────────────────────────────
% DOCUMENT
% ───────────────────────────────────────────────────────────────────────────
\begin{document}

% ═══════════════════════════════════════════════════════════════════════════
% HEADER
% ═══════════════════════════════════════════════════════════════════════════
\noindent
\begin{tikzpicture}[remember picture, overlay]
    \fill[fadarkblue] (current page.north west) rectangle ([yshift=-35mm]current page.north east);
\end{tikzpicture}

\vspace{-8mm}
\noindent
{\color{white}\sffamily\bfseries\Large EXECUTIVE SUMMARY}

\vspace{2mm}
\noindent
{\color{white}\sffamily\huge US-Iran Crisis Escalation Model}

\vspace{2mm}
\noindent
{\color{falightblue}\sffamily 28. Januar 2026 | EBF-S-2026-01-28-GEO-001}

\vspace{12mm}

% ═══════════════════════════════════════════════════════════════════════════
% MAIN CONTENT
% ═══════════════════════════════════════════════════════════════════════════

\sectionbox{FRAGESTELLUNG}{
Wie wahrscheinlich ist ein US-Militärschlag gegen Iran in den nächsten 48 Stunden?
}

\sectionbox{METHODIK}{
\textbf{Sequential Game} (Spieltheorie) mit 2 Hauptakteuren und je 3 Strategien, gelöst via \textbf{Backward Induction}. Parameter aus LLMMC-Prior + Bayesian Update mit aktueller Nachrichtenlage.

\textbf{Kontext-Dimensionen:} Temporal (frühe Amtszeit), Sozial (5'002 Tote bei Protesten), Ressourcen (USS Abraham Lincoln vor Ort), Kultur (Face-Saving beidseitig).
}

% ───────────────────────────────────────────────────────────────────────────
% KEY FINDING BOX
% ───────────────────────────────────────────────────────────────────────────
\highlightbox{
\begin{center}
{\color{fadarkblue}\sffamily\bfseries\Large KEY FINDING}

\vspace{3mm}
{\Huge\bfseries\color{fadarkblue} P(US-Strike in 48h) = 18\%}

\vspace{2mm}
{\small Konfidenzintervall: [12\% -- 25\%]}

\vspace{3mm}
\begin{tikzpicture}
    \fill[falightgray] (0,0) rectangle (12,0.5);
    \fill[fadarkblue] (0,0) rectangle (2.16,0.5);
    \node[right] at (2.3,0.25) {\small 18\%};
\end{tikzpicture}

\vspace{3mm}
{\color{fadarkgray}\textbf{Gleichgewicht:} Verhandlung $\rightarrow$ Deal (Pareto-optimal für beide Seiten)}
\end{center}
}

\sectionbox{GAME TREE (Vereinfacht)}{
\begin{center}
\begin{tikzpicture}[
    level 1/.style={sibling distance=40mm, level distance=12mm},
    level 2/.style={sibling distance=15mm, level distance=10mm},
    every node/.style={font=\small\sffamily},
    edge from parent/.style={draw, -latex}
]
\node[fill=fadarkblue, text=white, rounded corners, minimum width=15mm] {USA}
    child {node[fill=falightgray, rounded corners] {Drohen}
        child {node[font=\tiny] {Trotz}}
        child {node[font=\tiny] {Nachg.}}
    }
    child {node[fill=falightgray, rounded corners] {Strike}
        child {node[font=\tiny] {Vergelt.}}
        child {node[font=\tiny] {Absorb.}}
    }
    child {node[fill=fagreen!70, rounded corners, thick] {\textbf{Verhandeln}}
        child {node[fill=fagreen!30, font=\tiny, rounded corners] {\textbf{Accept}}}
        child {node[font=\tiny] {Reject}}
    };
\end{tikzpicture}

\vspace{2mm}
{\small\color{fadarkgray} Grün markiert: Gleichgewichtspfad (beste Option für beide)}
\end{center}
}

\sectionbox{TRIGGER-SZENARIEN}{
\begin{tabularx}{\textwidth}{@{}l X r r@{}}
\toprule
& \textbf{Trigger-Ereignis} & \textbf{P(Strike)} & \textbf{$\Delta$} \\
\midrule
\textcolor{fared}{\textbullet} & Iran exekutiert Demonstranten & 45\% & +27\% \\
\textcolor{fared}{\textbullet} & Proxy greift US-Truppen an & 60\% & +42\% \\
\textcolor{faorange}{\textbullet} & Trump-Tweet mit Deadline & 33\% & +15\% \\
\textcolor{fagreen}{\textbullet} & Direkter Trump-Iran Call & 5\% & -13\% \\
\bottomrule
\end{tabularx}

\vspace{2mm}
{\small\color{fadarkgray}\textbf{Aktuell:} Davos-Signal (24.01.) de-eskalierend, aber Armada vor Ort hält Drohkulisse aufrecht.}
}

\sectionbox{INTERPRETATION}{
«Das rationale Gleichgewicht ist Verhandlung -- beide Seiten gewinnen mehr durch einen Deal als durch Krieg. Die Gefahr liegt nicht im Kalkül, sondern in \textbf{Trigger-Ereignissen} (Hinrichtungen, Proxy-Angriffe), \textbf{Trumps Impulsivität} und \textbf{Irans Face-Saving-Zwang}.

Die maximale Drohkulisse dient der Verhandlung, nicht dem Krieg. Klassische \textbf{Coercive Diplomacy}.»
}

% ───────────────────────────────────────────────────────────────────────────
% FOOTER INFO
% ───────────────────────────────────────────────────────────────────────────
\vfill

\noindent\rule{\textwidth}{0.5pt}

{\small\color{fadarkgray}
\textbf{Quellen:} Al Jazeera, Military.com, Washington Post, RFE/RL (Januar 2026)

\textbf{Methodik:} Evidence-Based Framework (EBF), Game Theory (Selten 1965), LLMMC

\textbf{Modell-ID:} MOD-GEO-001 | \textbf{Session:} EBF-S-2026-01-28-GEO-001

\vspace{2mm}
\textbf{FehrAdvice \& Partners AG} | Prof. Ernst Fehr | www.fehradvice.com
}

\end{document}
